\documentclass[12pt,a4paper]{article}

% Language and encoding
\usepackage[utf8]{inputenc}
\usepackage[T1]{fontenc}
\usepackage[english]{babel}

% Mathematics and layout
\usepackage{amsmath,amssymb,amsthm}
\usepackage[a4paper,margin=2.5cm]{geometry}
\usepackage{xcolor}
\usepackage{listings}
\usepackage{hyperref}

% Hyperlinks
\hypersetup{
    colorlinks=true,
    linkcolor=blue,
    urlcolor=blue
}

% Python code formatting
\lstdefinestyle{pythonstyle}{
    language=Python,
    basicstyle=\ttfamily\small,
    keywordstyle=\color{blue},
    commentstyle=\color{gray},
    stringstyle=\color{red!70!black},
    numbers=left,
    numberstyle=\tiny\color{gray},
    stepnumber=1,
    numbersep=8pt,
    showstringspaces=false,
    breaklines=true,
    frame=single,
    tabsize=4
}

\lstset{style=pythonstyle}

% Mathematical environments
\newtheorem{axiom}{Axiom}
\newtheorem{theorem}{Theorem}

\title{\textbf{Alzheimer Maths}\\
\large The Mathematics of Alzheimer}
\author{Yami}
\date{}

\begin{document}

\maketitle

\begin{abstract}
This paper presents the foundations, theorems, and a computer
simulation of the system called \textit{Alzheimer Maths}.
\end{abstract}

\section{Foundations of the System: Axioms}

\subsection{The Initial Set}

The universe is initialized with a finite inventory containing the
integers from $1$ to $9$, their negatives from $-1$ to $-9$, their
reciprocals from $1$ to $\frac{1}{9}$, and their negative reciprocals.

The number $1$ appearing among the reciprocals is considered a
different iteration from the integer $1$, but behaves as $1$.

\subsection{Axiom 1: Destructive Consumption}

Every mathematical operation permanently consumes its operands.
A number that is used is destroyed and no longer exists, unless it
is recreated as the result of an operation.

The available operations are multiplication and addition. They
operate as they do within the ZFC system.

\subsection{Axiom 2: Multiplicity}

The same number may exist in several distinct iterations (tokens)
within the inventory.

\subsection{Axiom 3: Operation $P$ (Loss)}

The operation $P(x)$ consumes the token $x$ without generating
a new result. It is pure destruction.

\subsection{Axiom 4: Operation $R$ (Reset)}

The operation $R$ clears the current inventory and resets the system
exactly to its initial state.

\subsection{Axiom 5: No Reuse of a Single Token}

A term that exists in only one iteration cannot be used twice in
the same operation.

For example, if the set contains only the tokens $[1,0]$, the
operation
\[
1-1
\]
is forbidden because only one token with value $1$ is available.

\section{Fundamental Theorems}

\subsection{Theorem 1: The Number of Numbers}

\textit{``There will always be fewer numbers in the system than
there were at the beginning.''}

\paragraph{Proof.}
Each binary operation ($+$, $\times$, $-$, $\div$) consumes two
tokens and generates only one. The net change is therefore
$-1$ token.

Each operation $P$ consumes one token and generates none.
Its net change is also $-1$ token.

The total number of tokens therefore strictly decreases at every
step.

\subsection{Theorem 2: The End of the World}

\textit{``If you start with a set of $N$ numbers, you can perform
at most $N$ operations in this system, and the $N$-th operation
other than $R$ will always be a $P$.''}

\paragraph{Proof.}
According to Theorem 1, each operation removes exactly one token
from the system.

After $N-1$ operations, the inventory contains only one token:
\[
N-(N-1)=1.
\]

At this point, binary operations are impossible because there are
not enough operands. The only remaining legal action, other than
$R$, is the unary operation $P$.

The $N$-th operation that is not an $R$ must therefore be an
extinction through $P$.

\subsection{Theorem 3: The Maximum Possible Number}

Given that, for two positive integers $a$ and $b$ different from $1$,
\[
a \times b \geq a+b,
\]
the maximum possible number in our set is obtained by multiplying
the positive integers and then the appropriate negative factors.

At most eight negative factors whose absolute values exceed $1$
can be used:
\[
(-9)\times(-8)\times(-7)\times(-6)\times(-5)
\times(-4)\times(-3)\times(-2)=9!.
\]

We therefore obtain:
\[
(9!)^2.
\]

We can add to this number the sum of the positive fractions,
from $\frac{1}{1}$ to $\frac{1}{9}$.

Thus:
\[
(9!)^2
\left(
1+\frac12+\frac13+\frac14+\frac15+
\frac16+\frac17+\frac18+\frac19
\right)^2.
\]

For simplicity, let us call this number \textit{Max}:
\[
\boxed{
\mathrm{Max}=(9!)^2
\left(
1+\frac12+\frac13+\frac14+\frac15+
\frac16+\frac17+\frac18+\frac19
\right)^2
}
\]

Any other method is non-optimal: adding negative terms reduces the
value, multiplying fractions reduces the value, and adding positive
values produces numbers smaller than multiplying them.

This considers the set of possible strategies by examining each
group separately.

\paragraph{For the numbers from $1$ to $9$.}
They can be added or multiplied. Since
\[
a\times b\geq a+b,
\]
multiplying them produces a larger number.

\paragraph{For negative numbers.}
They can be added or multiplied. Adding them produces more negative
numbers. Even if we try to add them and then multiply by $-1$, we
will create a number smaller than a large product.

To create an even larger number using the first $9!$, the optimal
choice is to multiply them again.

\paragraph{For positive fractions.}
Since they lie between $-1$ and $1$, multiplying them always
produces smaller numbers. Adding them is the optimal choice.

Their sum will be greater than $1$, so multiplication is then
the optimal choice.

\paragraph{For negative fractions.}
The strategy follows the same logic as for the fractions: their sum
is used as a term in a product. To transform the negative sum into
a positive term in the product, we use the $-1$ left over from
the negative terms.

The resulting terms are greater than $1$. The optimal way to
combine them is therefore multiplication.

We obtain:
\[
\boxed{\mathrm{Max}=\ldots}
\]

\section{Python Simulation}

\subsection{Checking Token Multiplicity}

The \texttt{T2valide} function checks that an operation does not
reuse a single token twice.

\begin{lstlisting}
def T2valide(t1, t2, nombres):
    if t1 != t2:
        return True
    else:
        n = 0

        for i in nombres:
            if i == t1:
                n += 1

        return n > 1
\end{lstlisting}

\subsection{Simulation Program}

\begin{lstlisting}
print("Type stop to end the simulation")

Run = True
New_world = True


def T2valide(t1, t2, nombres):
    if t1 != t2:
        return True

    n = 0

    for i in nombres:
        if i == t1:
            n += 1

    return n > 1


while Run:

    if New_world:
        nombres = [
            # Positive integers
            1.0, 2.0, 3.0, 4.0, 5.0,
            6.0, 7.0, 8.0, 9.0,

            # Negative integers
            -1.0, -2.0, -3.0, -4.0, -5.0,
            -6.0, -7.0, -8.0, -9.0,

            # Positive fractions
            1.0, 1/2, 1/3, 1/4, 1/5,
            1/6, 1/7, 1/8, 1/9,

            # Negative fractions
            -1.0, -1/2, -1/3, -1/4, -1/5,
            -1/6, -1/7, -1/8, -1/9
        ]

        New_world = False
        print("New inventory:", nombres)

    else:
        operation = input(
            "Operation (+, *, P, R, stop): "
        )

        if operation == "stop":
            Run = False

        elif operation == "R":
            New_world = True

        elif operation == "P":
            T1 = float(input("Number to destroy: "))

            if T1 in nombres:
                nombres.remove(T1)
                print("Inventory:", nombres)
            else:
                print("Invalid operation: token not found.")

        elif operation in ("+", "*"):
            T1 = float(input("First number: "))

            if T1 not in nombres:
                print("First token not found.")
                continue

            T2 = float(input("Second number: "))

            if T2 not in nombres or not T2valide(
                T1, T2, nombres
            ):
                print("Invalid operation: tokens unavailable.")
                continue

            # Consume both tokens
            nombres.remove(T1)
            nombres.remove(T2)

            # Create the result
            if operation == "+":
                resultat = T1 + T2
            else:
                resultat = T1 * T2

            nombres.append(resultat)

            print("Result:", resultat)
            print("Inventory:", nombres)

        else:
            print("Unsupported operation.")
\end{lstlisting}

\end{document}
